\documentclass[10pt]{article}
\usepackage[margin=1in]{geometry}
\usepackage[T1]{fontenc}
\usepackage{lmodern}
\usepackage{xcolor}
\usepackage{listings}
\usepackage{hyperref}

\newcommand{\SONIC}{SONIC}
\newcommand{\SMITH}{\textsc{SMITH}}

\lstset{
  basicstyle=\ttfamily\scriptsize,
  breaklines=true,
  breakatwhitespace=false,
  columns=fullflexible,
  keepspaces=true,
  frame=single,
  rulecolor=\color{black!35},
  xleftmargin=0.5em,
  xrightmargin=0.5em
}

\title{Supporting Information for\\
SONIC: Symmetry-Oriented Non-redundant Internal Coordinates}
\author{Vincenzo Barone}
\date{}

\begin{document}
\maketitle

\section*{Norbornane SMITH/SONIC Standalone Input and Gaussian 16 Export}

This Supporting Information records the minimal standalone input required to
generate the norbornane \SONIC\ coordinate contract and the corresponding
Gaussian 16 \texttt{ReadAllGIC} input.  The first listing is intentionally not
the enriched \texttt{xyzin} container written after molecular perception and
coordinate construction.  It contains only a standard XYZ geometry followed by
the SMITH directives needed for this Gaussian-compatible export.

\subsection*{SMITH/SONIC input}

\lstinputlisting[caption={Minimal norbornane SMITH/SONIC standalone input: standard XYZ block plus the required export directives.}]{si/norbornane_smith_input.xyzin}

\subsection*{Generated Gaussian 16 input}

\lstinputlisting[caption={Gaussian 16 \texttt{ReadAllGIC} input generated automatically from the SMITH/SONIC contract.}]{si/norbornane_g16_readallgic.gjf}

\subsection*{Gaussian 16 output excerpt}

The complete Gaussian log is retained in the reproducibility archive.  The
excerpt below records the optimization progress and termination diagnostics.

\lstinputlisting[caption={Selected Gaussian 16 output lines for the norbornane \texttt{ReadAllGIC} run.}]{si/norbornane_g16_output_excerpt.txt}

\end{document}
